% Methodology paragraph for the paper

\paragraph{Methodology.}
All experiments follow a common evaluation protocol: stratified 10-fold cross-validation with an 80/20 train--test split per fold, using accuracy, macro F1, precision, recall, and ROC-AUC as evaluation metrics. Two embedding strategies are compared throughout: RoBERTa (\texttt{FacebookAI/roberta-base}), which produces 768-dimensional [CLS] embeddings, and Mistral (\texttt{mistral-embed}), which produces 1024-dimensional embeddings. Both kind of embeddings are precomputed before training. The two datasets differ in structure and therefore require distinct feature representations and model families.

For \textbf{EVONS}, articles are static documents with a title and a description. Each is encoded separately, and the two resulting 768-dimensional vectors are concatenated into a 1536-dimensional input. Three baseline classifiers are evaluated: logistic regression, random forest (250 trees), and XGBoost (250 estimators, maximum depth 6, learning rate 0.05). A two-hidden-layer MLP is also evaluated and constitutes the neural baseline. For virality prediction specifically, three additional MLP variants progressively incorporate non-textual signals: \textsc{SourceMLP} appends a learned 32-dimensional embedding of the publishing media source to the text representation; \textsc{AvgEngMLP} appends the standardised historical average engagement of that source; and \textsc{GatingMLP} replaces simple concatenation with a learned gating mechanism that softly interpolates between title and description embeddings conditioned on the engagement scalar, before passing the result to the classification head.

For \textbf{FakeNewsNet}, each sample is a Twitter propagation tree represented as an ordered sequence of tweets. Per-tweet RoBERTa or Mistral embeddings are computed individually, then averaged across the sequence to yield a single fixed-length vector (768- or 1024-dimensional respectively) that is fed to the same set of classifiers and the MLP described above. This mean-pooling strategy discards temporal ordering but provides a simple, sequence-length-agnostic representation compatible with standard classifiers.

Statistical comparisons between models use exact sign-flip permutation tests with bootstrap confidence intervals (10{,}000 resamples), with family-wise error control via Holm--Bonferroni correction over all 470 pairwise comparisons. Effect sizes are quantified with Cliff's $\delta$.